\section*{七、实验结论}
\addcontentsline{toc}{section}{七、实验结论}
本实验完成了基于 C++17 的 Linux 文件备份与恢复软件，实现目录扫描、规则过滤、多种归档与压缩算法、加密存储、完整及增量备份、定时与实时触发，以及基于 BPSEC1/BPSEC2 的远程备份。CLI 和 Qt 6 图形界面复用核心处理模块，服务端使用 SQLite 管理账户和快照信息。测试覆盖归档恢复、加密认证、增量依赖、图形交互、远端传输和软件安装。

系统围绕数据可恢复性组织处理流程：先验证材料，再在暂存目录重建文件；增量链按依赖顺序应用，全部成功后发布结果；保留策略保护仍被依赖的祖先快照。错误状态与诊断文本分离，临时文件创建、身份核对及清理相互配合，使中途失败能被识别并向用户说明。

\section*{八、总结及心得体会}
\addcontentsline{toc}{section}{八、总结及心得体会}
开发初期先完成目录树的备份与恢复，明确路径、文件类型和元数据表示，再将压缩、加密和增量策略接入字节流处理。CLI 与 GUI 复用同一核心，减少了过滤规则、编码选项和错误处理不一致的问题。通过 GNU tar、Zstd 等独立实现进行互操作测试，也帮助我们发现了仅靠自身编码和解码往返难以暴露的格式问题。

在增量恢复和临时文件清理的开发过程中，我们发现正常路径上的功能测试不足以覆盖中途失败。通过故障注入测试，进一步完善了错误状态向上层返回、临时资源管理和恢复结果发布的处理，也认识到备份软件必须重点保证失败时原有数据的完整性。

界面测试需要同时观察业务结果和用户看到的状态。实际桌面操作与自动化控件检查结合后，能够发现对话框取消、路径回填、任务忙碌及连接失败提示中的交互问题。小组成员在核心、界面、建模和文档工作中互相检查，使实现、图纸和测试能够相互对应。

\section*{九、对实验过程及方法手段的改进建议}
\addcontentsline{toc}{section}{九、对实验过程及方法手段的改进建议}
后续应继续完善文件系统异常处理，针对资源身份核验、目录替换和已有空目标的发布失败增加独立故障注入测试，检查错误返回、原有内容及临时文件处理结果。

性能测试可增加小文件密集目录、零填充大文件和随机数据，固定环境并重复测量编码、恢复与公网传输耗时，同时记录峰值内存和内容校验结果。自动任务和远端传输还可增加长时间运行与反复中断测试，检查事件积压、凭据失效和断点续传的状态恢复。
